\cotefile{67}
\scanpage{53}
% archive foliation: 55
\margin{\diagram{\begin{tikzcd}[ampersand replacement=\&,column sep=small]
\unsure{\text{torseurs sur groupes diédriques}} \arrow[r,leftrightarrow] \& \text{Cercles toriques}
\end{tikzcd}}
\deleted{\unsure{(syst. à torsion)}}
\deleted{$\unsure{\simeq\mathbb U\rtimes\mathbb Z/2\mathbb Z}$}
\deleted{\unsure{(comm.}\ $\unsure{\mathbb Z/2\mathbb Z}$\unsure{)}}}

\diagram{\begin{tikzcd}[ampersand replacement=\&,column sep=large,row sep=large]
\text{Cercles toriques} \arrow[r,leftrightarrow] \&
\begin{gathered}\deleted{\unsure{\text{plan}}}\ \text{vectoriel}\ \inserted{\text{réel}}\ \text{euclidien}\ \deleted{\unsure{\text{uni}}}\end{gathered}
\arrow[r,leftrightarrow] \&
\left[\begin{gathered}\unsure{\text{Droite vectorielle}}\ \inserted{\text{unitaire}}\ \text{sur corps }\mathbb C\\
\text{iso. à }\mathbb C\ \unsure{\text{(mais pas can. +)}}\end{gathered}\right]
\\
\text{Cercles toriques orientés}\arrow[r,leftrightarrow]\arrow[d,leftrightarrow] \&
\uline{\text{plans}}\ \text{vectoriel réel orienté}\arrow[r,leftrightarrow] \&
\begin{gathered}\deleted{\text{plan}}\ \inserted{\text{droites vectorielles}}\\
\text{complexes}\ \inserted{\text{unitaires}}\ \deleted{\unsure{D}}\ T\\
\deleted{\unsure{\text{structure euclidienne}}}\\
\text{i.e. avec iso }T\simeq\overline T^{\vee}\ \text{i.e.}\\ T^{\vee}\simeq\overline T\end{gathered}
\\
\text{torseurs sous }\mathbb U \& {}\arrow[r,leftrightarrow] \&
\begin{gathered}\text{droites vectorielles complexes}\\
\text{conformes, avec}\ \deleted{\unsure{\text{disques}}}\\
\deleted{\text{conformes}}\\
\unsure{\text{contour|contenant}}\ \unsure{\text{fixé}}\ D\end{gathered}
\end{tikzcd}}

\deleted{\unsure{Tandis que deux flèches de catégories}}

\deleted{Soient $C',C''$ deux cercles toriques orientés,}
\deleted{\unsure{alors on a}}
\deleted{$\unsure{\mathbb C\operatorname{Isom}}(C'^{-1},C'')\simeq
\operatorname{Isom}(C''^{-1},C')$}

\margin{\uline{NB} Pour $D$ droite conforme fixée, on a
\unsure{tous les}\ \unsure{jauges en}\ $D=\unsure{\delta\rtimes\mathbb U}$,
et $\unsure{\mathcal J_G(D)\simeq\mathcal J(D)^{\leftarrow\varepsilon}_{\frac12}}$ ;
(i.e. $\unsure{G}$ = groupe de \unsure{rotations}).}

\margin{\unsure{condition des}\ $\mathbb R^{*+}$\unsure{-torseurs}}

Le passage aux cercles toriques \inserted{orientés} duals correspond au passage
à l'espace vectoriel \deleted{réel euclidien}
\inserted{complexe dual} $\deleted{\unsure{T}}T^{\vee}$ $(\simeq\overline T)$.

La donnée, pour deux \uline{droites} vect complexes $T',T''$
d'un système $((D',D''),u)$, où $D',D''$ sont des
disques dans $T',T''$ et
\[
\begin{aligned}
u\in\delta D'\wedge_{\mathbb U}\delta D''\quad
(&\simeq\operatorname{Isom}_{\mathbb U}(\delta D'^{-1},\delta D'')\\
&\simeq\operatorname{Isom}_{\mathbb U}(\delta D''^{-1},\delta D')
\simeq\operatorname{Isom}(T'^{-1},T'')
\simeq\operatorname{Isom}(T''^{-1},T')
\simeq\unsure{\operatorname{Biisom}}(T'\otimes T'')\\
&\simeq\unsure{T'^{\times}\wedge_{\mathbb C^{\times}}T''^{\times}}
\unsure{\}}\simeq\deleted{\unsure{\text{flèches de }D',D''\ \text{pour}}}
\text{dualités entre }(T',T''))
\end{aligned}
\]
revient à la donnée d'une $v\in(T'\otimes_{\mathbb C}T'')^{\unsure{\times|*}}$
i.e. une dualité entre $T',T''$, \uline{plus} une jauge \inserted{$D'$} dans $T'$
(ou ce qui revient au \unsure{même}\ \unsure{quasi.}\ ; plus une jauge
$D''$ dans $\overline T''$).
